\begin{RequAnswer}{5.1}
$x_1=3-2=1$, $x_2=3^2-2^2=5$, $x_3=3^3-2^3=19$, $x_4=3^4-2^4=65$, and $x_5=3^5-2^5=211$.
\end{RequAnswer}
\begin{RequAnswer}{5.2}
It means to find an explicit formula (or closed formula) for it.
In other words, a formula that is not recursively defined.
\end{RequAnswer}
\begin{RequAnswer}{5.3}
$x_2=2 x_1+3 = 2\cdot 1+3=5$, $x_3=2 x_2+3 = 2\cdot 5+3=13$, $x_4=2 x_3+3=2\cdot 13+3=29$,
and $x_5=2 x_4+3=2\cdot 29+3=61$.
\end{RequAnswer}
\begin{RequAnswer}{5.4}
(a) $x_2=x_1+3 = 2+3=5$, $x_3=x_2+3 = 5+3=8$, $x_4=x_3+3=8+3=11$, and $x_5=x_4+3=11+3=14$.
(b) $x_n=3n-1$.
(c) Notice that if $x_n=3n-1$, then $x_{n-1}=3(n-1)-1=3n-4$. So
$x_{n-1}+3 = 3n-4+3=3n-1=x_n$, verifying that it works for the recurrence relation.
But we also need to show that it works for the base case: $x_1=2=3(1)-1$, so it also works for the base case and we
are done.
\end{RequAnswer}
\begin{RequAnswer}{5.5}
Increasing because for most algorithms, if you have more input it takes longer to even read the input, so it
would also take longer to run the algorithm. In rare cases it might be constant--for instance, an algorithm that
returns the fifth element of an array will take the same amount of time regardless of the size of the array.
\end{RequAnswer}
\begin{RequAnswer}{5.6}
They are sometimes monotonic. If $r<0$ they will not be monotonic, but if $r>0$ they are monotonic.
They are also sometimes increasing. For instance, if $r>1$, it will be increasing, but if
$0<r<1$ it will be decreasing. If $r<0$, they are neither increasing nor decreasing since they
go back and forth between positive and negative.
\end{RequAnswer}
\begin{RequAnswer}{5.7}
They are always monotonic. A given arithmetic progression is either always increasing or always decreasing,
depending on whether $d$ is positive or negative.
\end{RequAnswer}
\begin{RequAnswer}{5.8}
Many answers are possible, but your answer should be very similar in form as the ones given.
(a) $a_n=3(2/3)^n$.
(b) $b_n=2+8n$.
(c) $c_n=(5/3)c_{n-1}$, $c_0=3$.
(d) $d_n=d_{n-1}+9/4$, $d_0=8$.
\end{RequAnswer}
\begin{RequAnswer}{5.9}
They are not because $-x^i=-(x^i)$ and $(-x)^i=(-1)^i(x)^i$. Depending on whether $x$ is positive or negative
and depending on whether $i$ is even or odd, these may have opposite signs.
\end{RequAnswer}
\begin{RequAnswer}{5.10}
$\dsum_{i=0}^6 -(-3)^i$ or $\dsum_{i=0}^6 (-1)^{i+1}3^i$.
\end{RequAnswer}
\begin{RequAnswer}{5.11}
$\dsum_{k=0}^{30} 5k-7 = \dsum_{k=0}^{30} 5k - \dsum_{k=0}^{30} 7=5\dsum_{k=0}^{30} k - 7*31
=5\frac{30\cdot 31}{2} - 217 = 2108$
\end{RequAnswer}
\begin{RequAnswer}{5.12}
$\dsum_{k=0}^{n} 2^k = \frac{2^{n+1} -1}{2-1} = 2^{n+1}-1$
or
$\dsum_{k=0}^{n} 2^k = \frac{1 - 2^{n+1}}{1 - 2} = \frac{1-2^{n+1}}{-1} = 2^{n+1}-1$.
\end{RequAnswer}
\begin{RequAnswer}{5.13}
Sure. Whenever $x_0=0$.
\end{RequAnswer}
\begin{RequAnswer}{5.14}
If you can get this one without mistakes, you are doing {\em really} well! If you don't quite get it, keep
trying to do it on your own. You will learn a lot in the process and it will be a good algebra refresher.

$\displaystyle
\dsum_{k=1}^{23} \frac{11}{(-7)^k}
= 11\dsum_{k=1}^{23} \frac{1}{(-7)^k}
= 11\dsum_{k=1}^{23} \left(\frac{1}{-7}\right)^k
= 11\dsum_{k=1}^{23} \left(-\frac{1}{7}\right)^k
= 11\left(\dsum_{k=0}^{23} \left(-\frac{1}{7}\right)^k - \left(-\frac{1}{7}\right)^0\right)
$

$\displaystyle
= 11\left( \frac{1-\left(-\frac{1}{7}\right)^{24}}{1-\left(-\frac{1}{7}\right)} - 1\right)
= 11\left( \frac{1-\left(-1\right)^{24}\left(\frac{1}{7}\right)^{24}}{\frac{8}{7}} - 1\right)
= 11\left( \frac{7}{8}\left(1-\left(\frac{1}{7}\right)^{24}\right) - 1\right)
$

$\displaystyle
= 11\left(\frac{7}{8} -\frac{7}{8}\left(\frac{1}{7}\right)^{24} - 1\right)
= 11\left(-\frac{1}{8} -\frac{1}{8}\left(\frac{1}{7}\right)^{23}\right)
%= -\frac{11}{8} -\frac{11}{8}\left(\frac{1}{7}\right)^{23}
= -\frac{11}{8}\left(1+\left(\frac{1}{7}\right)^{23}\right)
%= \frac{77}{6}\left(1-\left(\frac{1}{7}\right)^{24}\right) - 11
%= \frac{11}{6}\left(\frac{1}{7}\right)^{23}+\frac{11}{6}
$.
\end{RequAnswer}
\begin{RequAnswer}{5.15}
According to Theorem~\ref{thm:maclaurin},
$\cos x = 1 - \dfrac{x^2}{2!} + \dfrac{x^4}{4!} - \cdots +(-1)^n\dfrac{x^{2n}}{(2n)!} +  \cdots$.
To approximate $\cos(1)$, we can just use several terms of the infinite sum.
So,
$$\cos(1) \approx 1 - \dfrac{1^2}{2!} + \dfrac{1^4}{4!} -  \dfrac{1^6}{6!}
=1-\dfrac{1}{2}+\dfrac{1}{24}-\dfrac{1}{720}
=\dfrac{720-360+30-1}{720} = \dfrac{389}{720}\approx 0.540277.$$
Note that $\cos(1)=0.540302305\cdots $, so our approximation is pretty good.
\end{RequAnswer}
\begin{RequAnswer}{5.16}
No. You can only add matrices that have the same
dimensions (number of rows and columns).
\end{RequAnswer}
\begin{RequAnswer}{5.17}
 $A = \begin{bmatrix}
1 & 3 & 9 \cr
2 & 6 & 18 \cr
4 & 12 & 36 \cr
8 & 24 & 72 \cr \end{bmatrix}.$
\end{RequAnswer}
\begin{RequAnswer}{5.18}
$3A=\begin{bmatrix}
-9& 0 & 3 \cr
12 & -6 & 15 \cr\end{bmatrix}$
$A-B=\begin{bmatrix}
-4& -7 & 6 \cr
-3 & -3 & -3 \cr\end{bmatrix}$.
\end{RequAnswer}
\begin{RequAnswer}{5.19}
${\bf 0}_{3\times 4}=
\begin{bmatrix}0 & 0 & 0 & 0 &0 \cr
0 & 0 & 0 & 0 &0\cr \end{bmatrix}$
and
 ${\bf I}_3=
\begin{bmatrix}
1 & 0 & 0 \cr
0 & 1 & 0 \cr
0 & 0 & 1 \cr\end{bmatrix}$.
\end{RequAnswer}
\begin{RequAnswer}{5.20}
In the first case you can since the number of columns of
the first matrix is equal to the number or rows of the second one.
The result will have dimension $3\times 7$.
In the second case you cannot do the multiplication since $3\neq 6$.
That is, the number of columns of the first is not the same
as the number or rows of the second.
\end{RequAnswer}
\begin{RequAnswer}{5.21}
$\begin{bmatrix} 2 & 2 \cr 0 & -2 \cr
\end{bmatrix}$
\end{RequAnswer}
\begin{RequAnswer}{5.22}
$A^{19}=A^5A^{14}={\bf 0}_n A^{14}={\bf 0}_n$.
\end{RequAnswer}
\begin{RequAnswer}{5.23}
 Observe that $$
\begin{bmatrix} -4 & x \cr -x & 4
\end{bmatrix}^2 =\begin{bmatrix} -4 & x \cr -x & 4
\end{bmatrix}\begin{bmatrix} -4 & x \cr -x & 4
\end{bmatrix} =  \begin{bmatrix} 16 - x^2 & 0 \cr 0 & 16 - x^2 \end{bmatrix}, $$ and so we must have
$16 - x^2 = -1$ which gives us $x = \pm\sqrt{17}$.
\end{RequAnswer}
\begin{RequAnswer}{5.24}
Disprove! Take for example $A = \begin{bmatrix} 0 & 0 \cr 1 & 1
\cr\end{bmatrix}$ and $B = \begin{bmatrix} 1 & 0 \cr 1 & 0
\cr\end{bmatrix}$. Then $$A^2 - B^2 = \begin{bmatrix} -1 & 0 \cr 0 &
1 \cr\end{bmatrix} \neq  \begin{bmatrix} -1 & 0 \cr -2 & 1
\cr\end{bmatrix}= (A+B)(A-B). $$
\end{RequAnswer}
\begin{RequAnswer}{5.25}
Disprove! This is not generally true. Take $A =
\dis{\begin{bmatrix} 1 & 1 \cr 1 & 2
\end{bmatrix}}$ and $B =
\dis{\begin{bmatrix} 3 & 0 \cr 0 & 1
\end{bmatrix}}$. Clearly $A^T = A$ and $B^T = B$. We have
$$AB  = \begin{bmatrix} 3 & 1 \cr 3 & 2\cr\end{bmatrix}$$but
$$(AB)^T  = \begin{bmatrix} 3 & 3 \cr 1 & 2\cr\end{bmatrix}.$$
\end{RequAnswer}
\begin{RequAnswer}{5.26}
Disprove! Take $A = B = {\bf I}_n$ for any $n>1$. Then $\tr{AB} = n< n^2 = \tr{A}\tr{B}$.  
\end{RequAnswer}
\begin{RequAnswer}{5.27}
 We have
$$(AA^T)^T = (A^T)^TA^T = AA^T.$$
\end{RequAnswer}
